% ============================================================
%  TEMA 2 — APARTADO 4: NUEVAS TENDENCIAS EN ADP — POR COMPETENCIAS
%
%  RECORDATORIO AL EQUIPO:
%  El Análisis y Diseño de Puestos (ADP) por competencias supone un cambio
%  de paradigma respecto al ADP tradicional:
%
%  ADP TRADICIONAL:
%    - Se centra en TAREAS y RESPONSABILIDADES del puesto.
%    - El resultado es una descripción de puesto estática.
%    - El encaje persona-puesto se hace comparando requisitos formales
%      (titulación, experiencia) con el perfil del candidato.
%
%  ADP POR COMPETENCIAS:
%    - Se centra en los COMPORTAMIENTOS OBSERVABLES que llevan al éxito.
%    - Identifica competencias clave (técnicas + conductuales/transversales).
%    - Es más flexible: la misma competencia puede ser válida en varios puestos.
%    - Facilita la gestión integrada del talento (selección, formación,
%      evaluación del desempeño, planes de carrera, retribución) bajo un
%      marco de competencias común.
%    - Permite adaptar rápidamente los perfiles a entornos cambiantes (VUCA).
%
%  OTRAS TENDENCIAS:
%    - ADP en entornos ágiles (puestos definidos por roles y sprints, no tareas).
%    - Uso de People Analytics e IA para rediseñar puestos basándose en datos.
%    - Gestión por proyectos vs. gestión por puestos (modelos gig/fluid).
%    - ADP continuo: revisión periódica sistemática en lugar de actualización
%      puntual y reactiva.
%    - Wellbeing y diseño de puestos orientado al bienestar del empleado.
%
%  PARA ESTA SECCIÓN INDICAD:
%  - ¿Aplica CaixaBank el ADP por competencias? ¿Tiene un diccionario de
%    competencias corporativo?
%  - ¿Cómo integra el enfoque de competencias en selección, formación y evaluación?
%  - ¿Usa People Analytics para analizar o rediseñar puestos?
%  - ¿Está la empresa satisfecha con su modelo? ¿Qué tendencias le faltan adoptar?
% ============================================================

\section{Nuevas tendencias en análisis y diseño de puestos: ADP por competencias}

\subsection{Del enfoque tradicional al enfoque por competencias}

El análisis y diseño de puestos (ADP) ha evolucionado desde el modelo
clásico, centrado en la descripción de tareas y responsabilidades, hacia
un enfoque basado en \textbf{competencias}, que pone el foco en los
comportamientos observables y medibles que conducen a un desempeño
excelente en el puesto.

\begin{table}[H]
  \centering
  \small
  \caption{Comparativa entre el ADP tradicional y el ADP por competencias}
  \begin{tabularx}{\textwidth}{p{0.23\textwidth} X X}
    \toprule
    \textbf{Dimensión}       & \textbf{ADP Tradicional} & \textbf{ADP por Competencias} \\
    \midrule
    Unidad de análisis       & Tareas y responsabilidades & Competencias y comportamientos \\
    Resultado                & Descripción de puesto estática & Perfil de competencias dinámico \\
    Orientación              & Presente                  & Presente y futuro \\
    Flexibilidad             & Baja                      & Alta \\
    Integración con otros procesos RRHH & Limitada   & Total (marco integrado) \\
    Adaptación al cambio     & Reactiva                  & Proactiva \\
    \bottomrule
  \end{tabularx}
\end{table}

\subsection{Aplicación del ADP por competencias en CaixaBank}

% TODO: CaixaBank, como gran empresa del IBEX 35, es probable que tenga
% un diccionario de competencias corporativo. Investigad:
%   - ¿Existe un modelo de competencias formalizado?
%   - ¿Cuáles son las competencias corporativas (transversales a todos los puestos)?
%   - ¿Hay competencias específicas por nivel o por familia de puestos?
%   - ¿Cómo se usa el modelo en selección, evaluación del desempeño y formación?

\subsubsection{Modelo de competencias de CaixaBank}

[Descripción a completar con la información de la entrevista y/o la
web corporativa de CaixaBank (sección de empleo / cultura)]

Ejemplos de competencias frecuentes en entidades bancarias (ajustar a CaixaBank):
\begin{itemize}
  \item \textbf{Orientación al cliente}: escucha activa, empatía, resolución
        de problemas del cliente.
  \item \textbf{Orientación a resultados}: planificación, seguimiento de objetivos,
        toma de decisiones basada en datos.
  \item \textbf{Trabajo en equipo y colaboración}: comunicación, confianza,
        gestión de conflictos.
  \item \textbf{Adaptabilidad e innovación}: apertura al cambio, aprendizaje
        continuo, pensamiento creativo.
  \item \textbf{Integridad y cumplimiento}: ética profesional, respeto normativo,
        confidencialidad.
\end{itemize}

\subsubsection{Integración del modelo de competencias en los procesos de RRHH}

% TODO: ¿Cómo se usa el diccionario de competencias en:
%   - Selección (entrevista por competencias, assesment centres)?
%   - Evaluación del desempeño (evaluación 360º, por objetivos + competencias)?
%   - Formación y desarrollo (planes de desarrollo individuales)?
%   - Planes de carrera y sucesión?

[Descripción a completar]

\subsection{Otras tendencias emergentes en ADP}

\subsubsection{People Analytics aplicado al análisis de puestos}

% TODO: ¿Usa CaixaBank datos para analizar patrones de desempeño, rotación
% o satisfacción por tipo de puesto? ¿Tienen un equipo de People Analytics?
% En una empresa de su tamaño es probable que sí. ¿Qué sistemas usan (SAP,
% Workday, Cornerstone...)?

[Descripción a completar]

\subsubsection{ADP en entornos ágiles}

% TODO: Especialmente relevante en el área de tecnología y digitalización
% de CaixaBank. ¿Trabajan con metodologías ágiles (Scrum, SAFe)?
% ¿Cómo afecta eso al diseño de puestos (roles vs. puestos fijos)?

[Descripción a completar]

\subsubsection{Diseño de puestos y bienestar del empleado (\textit{wellbeing})}

% TODO: CaixaBank ha publicado informes de sostenibilidad con compromisos
% en bienestar del empleado. ¿Se refleja eso en el diseño de los puestos?
% ¿Tienen indicadores de bienestar vinculados al diseño de puestos?

[Descripción a completar]

\begin{notaequipo}
  Preguntad en la entrevista:
  \begin{enumerate}
    \item ¿Tienen un diccionario o modelo de competencias corporativo?
          ¿Podéis verlo o que os lo describan?
    \item ¿Cómo se usa ese modelo en la selección y en la evaluación del desempeño?
    \item ¿Tienen un equipo o función de People Analytics?
    \item ¿Trabajan con metodologías ágiles en alguna área? ¿Cómo ha cambiado
          el concepto de ``puesto'' en esas áreas?
    \item ¿Qué tendencias en ADP creen que les falta adoptar?
  \end{enumerate}
\end{notaequipo}
